\documentclass[11pt,a4paper]{article}

% Packages
% XeLaTeX: use fontspec + a Unicode font so Vietnamese diacritics render correctly.
\usepackage{fontspec}
\usepackage{polyglossia}
\setmainlanguage{english}
\setotherlanguage{vietnamese}
\setmainfont{Times New Roman}
\setsansfont{Arial}
\setmonofont{Consolas}
\usepackage{geometry}
\usepackage{graphicx}
\usepackage{booktabs}
\usepackage{tabularx}
\usepackage{multirow}
\usepackage{amsmath}
\usepackage{amssymb}
\usepackage{hyperref}
\usepackage{xcolor}
\usepackage{float}
\usepackage{caption}
\usepackage{subcaption}
\usepackage{enumitem}
\usepackage{fancyhdr}
\usepackage{titlesec}
\usepackage{listings}
\usepackage{algorithm}
\usepackage{algpseudocode}

% Page setup
\geometry{margin=1in, headheight=14pt}
\setlength{\parindent}{0pt}
\setlength{\parskip}{6pt}

% Colors
\definecolor{primaryblue}{RGB}{0,84,159}
\definecolor{accentgreen}{RGB}{87,157,66}

% Hyperref setup
\hypersetup{
    colorlinks=true,
    linkcolor=black,
    urlcolor=primaryblue,
    citecolor=primaryblue
}

% Header/Footer
\pagestyle{fancy}
\fancyhf{}
\lhead{Seminar 2}
\rhead{Code Comment Classification}
\rfoot{Page \thepage}

% Title formatting
\titleformat{\section}{\large\bfseries\centering}{\thesection}{1em}{}
\titleformat{\subsection}{\normalsize\bfseries}{\thesubsection}{1em}{}

% New page before each section
\newcommand{\sectionbreak}{\clearpage}

% Code listing style
\lstset{
  basicstyle=\ttfamily\small,
  breaklines=true,
  frame=single,
  backgroundcolor=\color{gray!8},
  keywordstyle=\color{primaryblue},
  commentstyle=\color{accentgreen},
  stringstyle=\color{red!60!black},
  showstringspaces=false,
}

\begin{document}

% Title Page
\begin{titlepage}

\centering

% University or Institution header
{\large \textbf{HANOI UNIVERSITY OF SCIENCE AND TECHNOLOGY}}\\[0.3cm]
{\normalsize SCHOOL OF INFORMATION \& COMMUNICATION TECHNOLOGY}

\vspace{6cm}

% Main Title with enhanced typography
\begin{center}
{\Large\textbf{Code Comment Classification}}\\[0.3cm]
{\large Multi-label Classification of Python Code Comments}\\[0.2cm]
{\large (NLBSE'23 Shared Task --- Project 2)}
\end{center}

\vspace{2cm}

% Team information in a more professional layout
\begin{minipage}{0.8\textwidth}
\centering
\textbf{\large Team Members}

\vspace{0.5cm}

\begin{tabular}{@{}ll@{}}
\textbf{Nguyễn Khắc Thái Bình} & \texttt{20251324M} \\[0.3cm]
\textbf{Lã Minh Trung}         & \texttt{20251319M} \\[0.3cm]
\textbf{Nguyễn Thu Uyên}       & \texttt{20252279M} \\[0.3cm]
\textbf{Nguyễn Thị Nhã Linh}   & \texttt{20261262M} \\[0.3cm]
\textbf{Nguyễn Đức Anh}        & \texttt{20252469M} \\[0.3cm]
\end{tabular}
\end{minipage}

\vspace{2.5cm}

% Course details in a box
\begin{minipage}{0.7\textwidth}
\centering
\fbox{\begin{minipage}{0.9\textwidth}
\centering
\vspace{0.3cm}
\textbf{Course:} Seminar 2\\
\textbf{Instructor:} Dr. Phuong Nguyen (University of L'Aquila, Italy)\\
\textbf{Academic Year:} 2026-2027\\
\vspace{0.3cm}
\end{minipage}}
\end{minipage}

\vfill

{\large \textbf{\today}}

\end{titlepage}

\setcounter{tocdepth}{2}
\tableofcontents
\newpage

%==============================================================================
\section{Introduction}
\label{sec:introduction}
%==============================================================================

Code comments are an essential part of software maintenance: they explain the
intent behind a piece of code, document its parameters, and warn future
developers about pitfalls. However, comments are often inconsistent, incomplete,
or missing entirely. Automatically classifying the \emph{purpose} of a comment
sentence is therefore a useful building block for documentation tooling, code
review assistance, and comment-quality analysis. The NLBSE'23 Code Comment
Classification shared task \cite{nlbse2023} frames this as a \textbf{multi-label
classification} problem: a single comment sentence can serve several purposes at
once, and the model must predict the full set of applicable categories.

In this work we study the Python subset of the NLBSE'23 dataset. Each comment
sentence is assigned one or more of five categories:

\begin{itemize}[nosep]
    \item \textbf{Summary} --- a brief description of what the code does.
    \item \textbf{Expand} --- additional explanation of why/how the code works.
    \item \textbf{Parameters} --- documentation of function arguments.
    \item \textbf{Usage} --- guidance on how to use the code.
    \item \textbf{DevelopmentNotes} --- miscellaneous notes for developers
          (TODOs, warnings, reminders).
\end{itemize}

We make the following contributions:

\begin{itemize}[nosep]
    \item We identify and correct a critical label-encoding pitfall in the raw
          data: the dataset is stored in an ``exploded'' long format where each
          sentence appears five times, and the true label is the
          \texttt{instance\_type} column rather than the \texttt{category}
          column. Several early pipelines in the repository used the wrong
          column and produced degenerate results.
    \item We build a leakage-free evaluation protocol based on 5-fold
          \texttt{MultilabelStratifiedKFold}, refitting the TF-IDF vectorizer
          inside each fold.
    \item We benchmark four classical multi-label models (One-vs-Rest Logistic
          Regression, Classifier Chain, Linear SVM, SGD) and scaffold a
          \texttt{roberta-base} transformer baseline with class weighting.
    \item We provide a detailed per-label analysis showing that performance is
          strongly correlated with label frequency.
\end{itemize}

The remainder of this report is organized as follows.
Section~\ref{sec:dataset} describes the dataset and its structure.
Section~\ref{sec:preprocessing} details the preprocessing pipeline and system
architecture. Section~\ref{sec:models} presents the models.
Section~\ref{sec:evaluation} describes the cross-validation and evaluation
protocol. Section~\ref{sec:results} reports results, and
Section~\ref{sec:analysis} analyzes them. Section~\ref{sec:packaging} covers
project packaging and reproducibility, and Section~\ref{sec:conclusion}
concludes.

%==============================================================================
\section{Dataset \& Data Understanding}
\label{sec:dataset}
%==============================================================================

\subsection{Task Definition}

The NLBSE'23 Code Comment Classification task is a \textbf{multi-label}
classification problem. Formally, given a comment sentence $x$, the model must
predict a binary vector $\mathbf{y} \in \{0,1\}^{C}$ where $C = 5$ is the number
of categories and $y_c = 1$ indicates that category $c$ applies to $x$. Unlike
single-label classification, the categories are not mutually exclusive: a
comment may simultaneously be a \textit{Summary} and a \textit{Usage} note.

\subsection{Dataset Source and Size}

The full competition dataset spans three programming languages (Java, Python,
Pharo) with 6{,}738 comment sentences in total. This project focuses on the
\textbf{Python subset}: 2{,}555 unique comment sentences with 5 categories,
averaging \textbf{1.12 labels per sentence}. The relatively low average label
count confirms that most comments carry a single dominant purpose, while a
minority carry multiple.

\subsection{Raw Data Format (the ``Exploded'' Layout)}

The raw file \texttt{data/python.csv} stores the data in a \emph{long
one-vs-rest} layout: every comment sentence appears \textbf{five times}, once per
category. Table~\ref{tab:rawschema} describes the columns.

\begin{table}[H]
\centering
\caption{Schema of the raw \texttt{python.csv} file.}
\label{tab:rawschema}
\begin{tabular}{ll}
\toprule
\textbf{Column} & \textbf{Meaning} \\
\midrule
\texttt{comment\_sentence\_id} & unique identifier of the comment sentence \\
\texttt{class}                & source class/file containing the comment \\
\texttt{comment\_sentence}     & the comment text \\
\texttt{category}              & which of the 5 slots this row represents \\
\texttt{instance\_type}        & \textbf{the true binary label} (1 = belongs, 0 = not) \\
\bottomrule
\end{tabular}
\end{table}

\paragraph{Critical pitfall.}
Because each sentence is duplicated across five rows, the \texttt{category}
column merely names the slot --- it is \textbf{not} the label. The true label is
\texttt{instance\_type}. Several early pipelines in this repository incorrectly
used \texttt{category} as the target, which duplicated every comment five times
with mostly-contradictory single-label targets and produced degenerate results
(micro-F1 $\approx$ 0.01--0.33). Section~\ref{sec:preprocessing} describes the
fix.

%==============================================================================
\section{Data Preprocessing \& Architecture}
\label{sec:preprocessing}
%==============================================================================

\subsection{Correct Multi-label Encoding}

The core preprocessing step pivots the exploded rows back to \textbf{one row per
\texttt{comment\_sentence\_id}}, building a multi-hot label vector from
\texttt{instance\_type}. The logic (see
\texttt{src/models/ml/data\_utils.py}) is shown in
Listing~\ref{lst:load}.

\begin{lstlisting}[language=Python,caption={Correct multi-label loader.},label={lst:load}]
def load_multilabel(path='python.csv'):
    df = pd.read_csv(path, dtype=str, encoding='utf-8', on_bad_lines='skip')
    df['instance_type'] = df['instance_type'].astype(int)
    categories = sorted(df['category'].unique())
    pivot = df.pivot_table(
        index=['comment_sentence_id', 'class', 'comment_sentence'],
        columns='category', values='instance_type',
        aggfunc='max', fill_value=0,
    ).reset_index()
    pivot = pivot.sort_values('comment_sentence_id').reset_index(drop=True)
    texts = pivot['comment_sentence'].fillna('').astype(str).values
    Y = pivot[categories].values.astype(int)
    return texts, Y, categories, pivot
\end{lstlisting}

The result is a clean multi-label matrix $Y \in \{0,1\}^{2555 \times 5}$ with no
missing values and no duplicate \texttt{comment\_sentence\_id}. The positive
count of each label is reconciled against the original file to guarantee that no
information is lost during the pivot.

\subsection{Text Normalization}

Comment text is normalized by collapsing whitespace and stripping leading and
trailing spaces. No lower-casing or stop-word removal is applied for the
transformer track; the TF-IDF track relies on n-grams to capture lexical
patterns, including case-sensitive identifiers.

\subsection{Feature Engineering (Classical Track)}

Two complementary TF-IDF views are combined with a \texttt{FeatureUnion}:

\begin{itemize}[nosep]
    \item \textbf{Word n-grams}: 1--2 grams, up to 20{,}000 features.
    \item \textbf{Character n-grams} (\texttt{char\_wb}): 3--5 grams, up to
          20{,}000 features.
\end{itemize}

Character n-grams help with code identifiers, camelCase tokens, and punctuation
patterns that word tokenization would split awkwardly. The two views are
concatenated into a single sparse feature matrix.

\subsection{Precomputed Artifacts}

The repository ships precomputed, version-controlled artifacts so that every team
member trains on identical splits. Table~\ref{tab:artifacts} summarizes them.

\begin{table}[H]
\centering
\caption{Precomputed artifacts committed to the repository.}
\label{tab:artifacts}
\begin{tabular}{lll}
\toprule
\textbf{Track} & \textbf{Directory} & \textbf{Contents} \\
\midrule
TF-IDF  & \texttt{data/tfidf/}   & \texttt{X\_train.npz}, \texttt{X\_test.npz}, vectorizers, folds \\
RoBERTa & \texttt{data/roberta/} & token ids, attention masks, labels, class weights, folds \\
\bottomrule
\end{tabular}
\end{table}

Shared fold definitions live in \texttt{folds.json} and are reused across both
tracks to guarantee comparability of results.

\subsection{System Architecture}

Figure~\ref{fig:architecture} illustrates the end-to-end pipeline. Raw
\texttt{python.csv} is pivoted into a multi-label table, then split into two
parallel tracks: a TF-IDF track feeding classical linear models, and a
tokenization track feeding the RoBERTa transformer. Both tracks consume the same
cross-validation folds and emit metrics into \texttt{results/}.

\begin{figure}[H]
    \centering
    \fbox{\begin{minipage}{0.9\textwidth}
        \centering
        \vspace{0.3cm}
        \texttt{python.csv} (exploded) \\
        $\downarrow$ pivot on \texttt{instance\_type} \\
        multi-label table (2{,}555 $\times$ 5) \\
        $\swarrow$ \hspace{3cm} $\searrow$ \\
        TF-IDF (word + char) \hspace{1.5cm} RoBERTa tokenizer \\
        $\downarrow$ \hspace{4cm} $\downarrow$ \\
        Linear models \hspace{2.5cm} RoBERTa + sigmoid head \\
        $\searrow$ \hspace{3cm} $\swarrow$ \\
        shared 5-fold CV $\rightarrow$ \texttt{results/}
        \vspace{0.3cm}
    \end{minipage}}
    \caption{End-to-end system architecture. Both modelling tracks share the
    corrected multi-label table and the same cross-validation folds.}
    \label{fig:architecture}
\end{figure}

%==============================================================================
\section{Model Development}
\label{sec:models}
%==============================================================================

\subsection{Classical Baselines}

Four linear models are trained on the TF-IDF features:

\begin{enumerate}[nosep]
    \item \textbf{One-vs-Rest Logistic Regression} (best overall).
    \item \textbf{Classifier Chain} with Logistic Regression base estimator.
    \item \textbf{One-vs-Rest Linear SVM}.
    \item \textbf{One-vs-Rest SGD} with log-loss.
\end{enumerate}

The One-vs-Rest strategy trains one independent binary classifier per label,
which is simple and parallelizable. The Classifier Chain instead trains labels
sequentially, feeding each label's prediction as an additional feature to the
next, thereby modelling label co-occurrence.

\paragraph{Leakage control.}
The vectorizer is \textbf{refit inside each cross-validation fold} on the
training partition only, so validation folds never influence the feature space.
This is a deliberate correction over earlier pipelines that fit the vectorizer
globally and therefore leaked information from the validation set.

\subsection{RoBERTa Baseline}

A \texttt{roberta-base} model is fine-tuned with a multi-label head:

\begin{itemize}[nosep]
    \item CLS representation (768-dim) $\rightarrow$
          \texttt{Linear(768, num\_labels)}.
    \item Loss: \texttt{BCEWithLogitsLoss} with per-class positive weights
          (\texttt{class\_weights.json}) to counter label imbalance.
    \item Max sequence length 128, learning rate $2\times10^{-5}$, weight decay
          0.01, warmup ratio 0.1, early-stopping patience 2.
    \item Optional per-category label masking for category-specific splits.
\end{itemize}

The model definition lives in \texttt{src/models/dl/roberta\_baseline.py} and the
training loop in \texttt{experiments/dl/train\_exp1.py}. Configuration is stored
in \texttt{configs/exp1\_roberta\_baseline.yaml}.

%==============================================================================
\section{Cross-validation \& Evaluation}
\label{sec:evaluation}
%==============================================================================

\subsection{Validation Protocol}

We use \textbf{5-fold \texttt{MultilabelStratifiedKFold}} (from the
\texttt{iterative-stratification} package), which preserves the joint label
distribution across folds. This is essential for multi-label data: naive KFold
can produce folds that entirely miss rare labels such as
\textit{DevelopmentNotes}, making the evaluation unstable.

\subsection{Metrics}

\begin{itemize}[nosep]
    \item \textbf{micro-F1} --- aggregates TP/FP/FN globally; dominated by
          frequent labels.
    \item \textbf{macro-F1} --- unweighted mean over labels; sensitive to rare
          labels.
    \item \textbf{Jaccard (samples)} --- exact-set overlap, rewards correct
          label combinations.
    \item \textbf{Hamming loss} --- fraction of misclassified label slots.
    \item \textbf{micro precision / recall}.
\end{itemize}

%==============================================================================
\section{Results}
\label{sec:results}
%==============================================================================

\subsection{Classical Models (5-fold CV, mean)}

Table~\ref{tab:classical} reports the mean cross-validation results for the four
classical models.

\begin{table}[H]
\centering
\caption{Classical ML results on the corrected multi-label Python subset.}
\label{tab:classical}
\begin{tabular}{lccccc}
\toprule
\textbf{Model} & \textbf{micro-F1} & \textbf{macro-F1} & \textbf{Jaccard} & \textbf{Hamming} & \textbf{Prec/Rec} \\
\midrule
\textbf{OVR Logistic Regression} & \textbf{0.6282} & \textbf{0.5928} & 0.5843 & 0.1697 & 0.6172 / 0.6397 \\
OVR SGD (log-loss)               & 0.6266 & 0.5932 & 0.5876 & 0.1746 & 0.6018 / 0.6536 \\
OVR Linear SVM                   & 0.6236 & 0.5854 & 0.5801 & 0.1709 & 0.6160 / 0.6313 \\
Classifier Chain (LR)            & 0.6098 & 0.5667 & \textbf{0.5977} & 0.1748 & 0.6103 / 0.6093 \\
\bottomrule
\end{tabular}
\end{table}

\subsection{Per-label Breakdown (OVR Logistic Regression)}

Table~\ref{tab:perlabel} breaks down the best model's performance by category.

\begin{table}[H]
\centering
\caption{Per-label precision/recall/F1 for the best model.}
\label{tab:perlabel}
\begin{tabular}{lcccc}
\toprule
\textbf{Label} & \textbf{Precision} & \textbf{Recall} & \textbf{F1} & \textbf{Support} \\
\midrule
Usage            & 0.7251 & 0.7000 & 0.7118 & 800 \\
Parameters       & 0.7035 & 0.7090 & 0.7050 & 794 \\
Expand           & 0.5620 & 0.6032 & 0.5815 & 504 \\
Summary          & 0.5551 & 0.6079 & 0.5793 & 454 \\
DevelopmentNotes & 0.3639 & 0.4138 & 0.3862 & 312 \\
\bottomrule
\end{tabular}
\end{table}

\subsection{RoBERTa Baseline}

\textit{Status: to be completed.} The RoBERTa track is scaffolded and its
preprocessed tensors are committed under \texttt{data/roberta/}. Training must be
run on a machine with sufficient disk space for the $\sim$500\,MB
\texttt{roberta-base} checkpoint. Results will be added here once available.

%==============================================================================
\section{Result Analysis}
\label{sec:analysis}
%==============================================================================

\subsection{Key Observations}

\begin{itemize}
    \item \textbf{Performance tracks label frequency.} \textit{Usage} and
          \textit{Parameters} (support $\approx$ 800) reach F1 $\approx$
          0.70--0.71, while \textit{DevelopmentNotes} (support 312) lags at F1
          $\approx$ 0.39. More training examples per label directly improve that
          label's F1.

    \item \textbf{Linear models are near-equivalent.} LR, SVM, and SGD all land
          in the 0.61--0.63 micro-F1 band, suggesting the TF-IDF space is close
          to linearly separable and the bottleneck is data volume and label
          imbalance rather than model capacity.

    \item \textbf{Classifier Chain trades F1 for Jaccard.} It achieves the best
          Jaccard (0.5977) by modelling label co-occurrence (a comment tagged
          \textit{Usage} is likelier to also carry \textit{Parameters}), at a
          small cost to per-label F1.

    \item \textbf{DevelopmentNotes is intrinsically hard.} It is both the rarest
          label and the most semantically diffuse (``miscellaneous notes''),
          explaining its clear margin of under-performance.
\end{itemize}

\subsection{Error Analysis}

\textit{To be completed.} Planned analyses include confusion patterns between
\textit{Summary} and \textit{Expand}, sensitivity to the decision threshold, and
qualitative examples of misclassified comments.

%==============================================================================
\section{Project Packaging \& Reproducibility}
\label{sec:packaging}
%==============================================================================

\subsection{Repository Layout}

\begin{lstlisting}
Seminar2/
  configs/            # experiment YAML configs
  data/               # raw + preprocessed data (tfidf/, roberta/)
  docs/report/        # this report
  experiments/        # dl/ and ml/ experiment entry points
  notebooks/          # data preprocessing + EDA notebooks
  results/            # generated metrics and analysis
  src/
    data/             # dataset loaders (nlbse2023.py)
    models/ml/        # classical pipelines
    models/dl/        # RoBERTa model definition
\end{lstlisting}

\subsection{Reproduction Steps}

\begin{enumerate}[nosep]
    \item Install dependencies: \texttt{pip install -r requirements.txt}.
    \item (Optional) Regenerate preprocessed artifacts from the notebooks in
          \texttt{notebooks/}.
    \item Run classical baselines:
          \texttt{python -m src.models.ml.ml\_pipeline}.
    \item Run RoBERTa baseline:
          \texttt{python -m experiments.dl.train\_exp1}.
    \item Inspect metrics under \texttt{results/}.
\end{enumerate}

\subsection{Known Limitations}

\begin{itemize}[nosep]
    \item Python-only subset (2{,}555 samples), not the full 6{,}738-sample
          Java+Python+Pharo dataset; each language has a different label set.
    \item RoBERTa results pending (disk-space constraint).
    \item Fixed 0.5 decision threshold for OVR/SGD; no per-label threshold
          tuning was performed.
\end{itemize}

%==============================================================================
\section{Discussion, Conclusion, and Future Work}
\label{sec:conclusion}
%==============================================================================

\subsection{Discussion}

The corrected preprocessing pipeline transformed a degenerate baseline
(micro-F1 $\approx$ 0.01) into a strong classical result (micro-F1 0.6282). This
underscores that in multi-label tasks, getting the label encoding right matters
far more than model choice: all four linear models cluster tightly around the
same performance, and the dominant source of error is label imbalance rather
than insufficient model capacity.

\subsection{Conclusion}

We corrected a fundamental label-encoding error in the raw data, established a
leakage-free 5-fold cross-validation protocol, and benchmarked four classical
multi-label models alongside a scaffolded RoBERTa baseline. The best classical
model achieves micro-F1 0.6282 and macro-F1 0.5928, with performance strongly
correlated to per-label support.

\subsection{Future Work}

\begin{itemize}[nosep]
    \item \textbf{Complete the RoBERTa comparison}: run the transformer baseline
          and compare it against the classical models.
    \item \textbf{Per-label threshold tuning}: optimize the decision threshold
          for each label independently to improve rare-label F1.
    \item \textbf{Extend to the full multilingual dataset}: incorporate the Java
          and Pharo subsets for a complete NLBSE'23 evaluation.
    \item \textbf{Error analysis}: systematically characterize the
          \textit{Summary}/\textit{Expand} confusion and other failure modes.
\end{itemize}

%==============================================================================
% References
%==============================================================================

\addcontentsline{toc}{section}{References}

\begin{thebibliography}{9}

\bibitem{nlbse2023}
R.~Karanikolas, C.~Aivaloglou, M.~Wessel, A.~Serebrenik, and A.~Bacchelli,
``NLBSE'23 tool competition: Code comment classification,''
in \textit{Proc.\ IEEE/ACM Int.\ Workshop on Natural Language-Based Software
Engineering (NLBSE)}, 2023.

\bibitem{pedregosa2011sklearn}
F.~Pedregosa \textit{et al.},
``Scikit-learn: Machine learning in Python,''
\textit{Journal of Machine Learning Research}, vol.~12, pp.~2825--2830, 2011.

\bibitem{sechidis2011stratification}
K.~Sechidis, G.~Tsoumakas, and I.~Vlahavas,
``On the stratification of multi-label data,''
in \textit{Proc.\ ECML PKDD}, 2011.

\bibitem{liu2019roberta}
Y.~Liu \textit{et al.},
``RoBERTa: A robustly optimized BERT pretraining approach,''
arXiv:1907.11692, 2019.

\bibitem{zhang2018mlknn}
M.-L. Zhang and Z.-H. Zhou,
``A review on multi-label learning algorithms,''
\textit{IEEE Trans.\ Knowledge and Data Engineering}, vol.~26, no.~8,
pp.~1819--1837, 2014.

\end{thebibliography}

\end{document}
